% The next line makes it so it compiles the main file if I hit the compile button
% when editing this file in TeXworks.
% !TEX root = ../AIDMA.tex
\section{Problems}

\begin{rem}
For the remainder of the book, whenever a problem asks for an algorithm, 
always assume it is asking for the most efficient
algorithm you can find.  You will likely lose points if your algorithm is not efficient enough.
We will use our intuitive definition of {\em efficient} until
we study algorithm analysis in 
Chapter~\ref{ch:AlgAn}.
\end{rem}

\begin{pro}
You are helping a friend debug the code below.  
He tells you ``The code in the if statement never executes.  
I have tried it for \verb+x=2+, \verb+x=4+, and even \verb+x=-1+, 
and it never gets to the code inside the if statement.''
%The following comments should not be needed in this context.  CAC, 7/30/13
%(Note: The statement \verb+x%2==0+ 
%is checking the parity of \verb+x+.  It is true if \verb+x+ is even and false if \verb+x+ is odd.
%Also, \verb+&&+ is the logical AND operator and \verb+||+ is the logical OR operator.)
\begin{code}
    if((x%2==0 && x<0) || !(x%2==0 || x<0)) {
        // Do something.
    }
\end{code}
\begin{enumerate}[(a)]
\item
Is he correct that the code inside the if statement does not execute for his chosen values?  
Justify your answer.
\item
Under what conditions, if any, will the code in the if statement execute?  Be specific and complete.
\item
Rewrite the conditional statement so that it is simpler (i.e. easier to understand what the condition is).
\end{enumerate}
\end{pro}

\begin{pro}
Simplify the following code as much as possible:
\begin{code}
	if(x<=0 && x>0) {
	    doSomething();
	} else {
	    doAnotherThing();
	}
	
\end{code}
\end{pro}

\begin{pro}
Simplify the following code as much as possible. 
(It can be simplified into a single if statement that is about as complex as the original outer if statement).
\begin{code} 
    if ( (!x.size()<=0  && x.get(0)!=11) || x.size()>0 ) {
       if ( !(x.get(0)==11 && (x.size()>13 || x.size()<13) ) 
          && (x.size()>0 || x.size()==13) ) {
                // Do a few things.
       }
    }
\end{code}
\end{pro}

% Reordering means this is no longer relevant.
%\begin{pro}\label{modEquivRel}
%Let $n$ be a positive integer.
%Recall that $a\equiv b\pmod n$ iff $n$ divides $a-b$ (that is,
%$a-b=k\cdot n$ for some $k\in\integers$).
%Use this formal definition to prove each of the following:
%\begin{lenum}
%\item
%$a\equiv a \pmod n$. (Reflexive property)
%\item
%If $a\equiv b\pmod n$, then $b\equiv a \pmod n$. (Symmetric property)
%\item
%If $a\equiv b \pmod n$ and $b\equiv c \pmod n$, then $a\equiv c \pmod n$. (Transitive property)
%\end{lenum}
%\end{pro}


\begin{pro}
Implement the swap operation for integers without using an additional variable
and without using addition or subtraction. (Hint: bit operations)
\end{pro}

\begin{pro}
Prove or disprove that the following method correctly computes the maximum of two integers \verb+x+ and \verb+y+,
assuming that the \verb+minimum+ method correctly computes the minimum of \verb+x+ and \verb+y+.
\begin{code}
    int maximum(int x, int y) {
        int min = minimum(x,y);
        int max = x + y - min;
        return max;
    }
\end{code}
\end{pro}


\begin{pro}
Give an algorithm \verb+int sumLarge(int []a,int n,int k)+
(where $a$ is an array with $n$ elements) 
 that returns the sum of all of the elements 
from the array $a$ that are at least $k$ (as in $a[i]\geq k$, not whose index is at least $k$).
\end{pro}


\begin{pro}
Let $a$ be an array with $n$ elements.
Give an algorithm \verb+int sum(int []a,int n)+ that returns the sum of all of the elements 
from the array $a$.
\end{pro}


\begin{pro}
Let $a$ be an array with $n$ elements.
Give an algorithm that returns true if and only if the elements of $a$ are in increasing order.
That is, if $i<j$, then $a[i]\leq a[j]$.
Call your algorithm \verb+boolean increasing(int []a,int n)+.
(Hint: You do not have to check that $a[i]\leq a[j]$ for every $i<j$. What is a simpler
thing to check that is equivalent to this?)
\end{pro}

\begin{pro}
Let $a$ be an array with $n$ elements.
Give an algorithm \verb+sort(int []a,int n)+ 
that sorts the elements of the array $a$ in increasing order.
\end{pro}


\begin{pro}
Let $a$ be an array with $n$ elements and assume that $a[i]$ is the number of 
people in room $i$ of a hotel.
The hotel can have at most \verb+capacity+ total guests. If they have too many guests,
they have to kick guests out in reverse order of room number (so larger room numbers get kicked out first).
Give an algorithm \verb+int tooMany(int []a,int n, int capacity)+ 
that returns the room number $k$ such that some or all guests in room $k$ and all guests 
in rooms $k+1$ and beyond must vacate;
or -1 if there is enough space for everybody. 
\end{pro}


\begin{pro}
Give a recursive algorithm that computes $n!$.  You can assume $n\geq 0$.
\end{pro}

\begin{pro}
Let $a$ be an array with $n$ elements.
Give an algorithm \verb+mod(int []a,int n,int k)+ that 
replaces each element $a[i]$ from the array $a$ with $a[i]\bmod k$. 
\end{pro}
\begin{pro}
Consider the following code.
{\small
\begin{code}
	boolean notBothZero(int x, int y) {
	    if(!(x==0 && y==0)) {
	        return true;
	    } else {
	        return false;
	    }
	}
	
	boolean unknown1(int x, int y) {
	    if(x!=0 && y!=0) {
	        return true;
	    } else {
	        return false;
	    }
	}
	
	boolean unknown2(int x, int y) {
	    if(x!=0 || y!=0) {
	        return true;
	    } else {
	        return false;
	    }
	}
\end{code}
}
\begin{lenum}
\item
Is \verb+unknown1+ equivalent to \verb+notBothZero+? Prove or disprove it.
\item
Is \verb+unknown2+ equivalent to \verb+notBothZero+? Prove or disprove it.
\item
Are \verb+unknown1+  and \verb+unknown2+ equivalent to each other? Prove or disprove it.
\end{lenum}
\end{pro}


\begin{pro}
The following method returns true if and only if none of the entries of the array are 0:
\begin{code}
	boolean noZeroElements(int[] a, int n) {
	    for(int i=0;i<n;i++) {
	        if(a[i] == 0 )
	            return false;
	    }
	    return true;
	}
\end{code}

The two methods below implement this idea for two arrays.
Assume \verb+list1+ and \verb+list2+ have the same size for both of these methods.
{\small
\begin{code}
	boolean unknown1(int[] list1, int[] list2, int n) {
	    for(int i=0;i<n;i++) {
	        if( list1[i]==0 && list2[i]==0 )
	            return false;
	    }
	    return true;
	}
	
	boolean unknown2(int[] list1, int[] list2, int n) {
	    if(noZeroElements(list1, n)) {
	        return true;
	    } else if(noZeroElements(list2, n) {
	        return true;
	    } else {
	        return false;
	    }
	}
\end{code}
}
\begin{enumerate}[(a)]
\item
What is \verb+unknown1+ determining? (Give answer in terms of \verb+list1+ and \verb+list2+ and the
appropriate quantifier(s).)
\item
What is \verb+unknown2+ determining? (Give answer in terms of \verb+list1+ and \verb+list2+ and the
appropriate quantifier(s).)
\item
Prove or disprove that \verb+unknown1+ and \verb+unknown2+ are determining the same thing.
\end{enumerate}
\end{pro}





\begin{pro}
Give an algorithm \verb+secondSmallest(int []a,int n)+
that returns the second smallest element of an array $a$ with $n$ elements.
Implement this under two different assumptions:
\begin{enumerate}[(a)]
\item
If there are two or more of the smallest value in the array, that value should be returned.
(e.g. If the array is $[2, 5, 6, 2, 4, 3, 2]$, then $2$ should be returned.)
\item
If there are two or more of the smallest value in the array, return the next larger value.
(e.g. If the array is $[2, 5, 6, 2, 4, 3, 2]$, then $3$ should be returned.)
\end{enumerate}
(Note: If there are not two or more of the smallest value, you of course return the next larger value
than the smallest value in either case.)
\end{pro}

\begin{pro}
Give an algorithm that will round a real number \verb+x+ to
the closest integer, rounding {\em up} at .5.
You can {\em only} use \verb+floor(y)+, \verb+ceiling(y)+, basic arithmetic (+, -, *, /) and/or numbers.
You {\em cannot} use anything else, including conditional statements!
{\em Prove that your algorithm is correct.}
\end{pro}

\begin{pro}
What will the following algorithm return for $n = 3$?
\begin{code}
	iCanDuzSomething(int n) {
	    int x = 0;
	    while(n>0) {
	        for(int i=1;i<=n;i++) {
	            for(int j=i;j<=n;j++) {
	                x = x + i*j;
	            } 
	        }
	        n--;
	    }
	    return x;
	}
\end{code}
\end{pro}


\begin{pro}
Recall that Example~\ref{exer:roundNotTruncate} had the conditions that $n>0$ and $m>2$.
Also recall that you gave a solution to this in Exercise~\ref{exer:roundNotTruncatetwo}.
Also recall that integer division always truncates {\em toward zero}, so negative numbers truncate
differently than positive ones.
\begin{lenum}
\item
Does your solution work when $m=2$?  Justify your answer with a proof/counterexample.
\item
Does your solution work when $n\leq 0$?  Justify your answer with a proof/counterexample.
\item
Give an algorithm that will work for any integer $n$ and any non-zero $m$.
Give examples that demonstrate that your algorithm is correct for the various cases and/or
a proof that it always works.  Make sure you consider all relevant cases (e.g., when it should
round up and down, when $n$ and $m$ are positive/negative).
You may only use basic integer arithmetic and conditional statements.
You may {\em not} use {\em floor}, {\em ceiling}, {\em abs} (absolute value), etc.  
You also may {\em not} use the {\em mod} operator since how it works with negative numbers 
is not the same for every language.
\end{lenum}
\end{pro}

\begin{pro}
Assume you have a function \verb+random(int n)+ that returns a random integer between $0$ and $n-1$, inclusive.
Give code/pseudocode for an algorithm \verb+random(int a, int b)+ that returns a random number between $a$ and $b$, inclusive of both $a$ and $b$.
You may assume that $a<b$ (although in practice, this should be checked).  
You may only call \verb+random(int n)+ once and you may not use conditional statements.
{\em Prove that your algorithm returns an integer in the required range.}
\end{pro}

\begin{pro}
Assume you have a function \verb+random()+ that returns a non-negative random integer.
Give code/pseudocode for an algorithm \verb+random(int a, int b)+ that returns a random integer 
between $a$ and $b$, inclusive of both $a$ and $b$.  
Each possible number generated should occur with approximately the same probability.
You may assume that $a$ and $b$ are both positive and that $a<b$ (although in practice, this should be checked).  
You may use only basic integer arithmetic (including the mod operator) and you may only call  \verb+random()+ once.
You may not use loops, conditional statements,  {\em floor}, {\em ceiling}, {\em abs} (absolute value), etc.
{\em Prove that your algorithm returns an integer in the required range.}
\end{pro}

\begin{pro}
Give an algorithm that prints
all of the primes that are less than or equal to $n$.  
Your algorithm should be as efficient as possible.
One approach is to modify the algorithm from Example~\ref{exa:isprime1} 
by using an array to make it more efficient.
\end{pro}

\begin{pro}
The following method is a simplified version of a method that might be used to 
implement a hash table or in a cryptographic system. 
Assume that for one particular use the number returned by this function
has to have the opposite parity (even/odd) of the parameter.  
For instance, \verb+hash_it(4)+ returns 49 which has the opposite 
parity of 4, so it works for 4.
Prove or disprove that this function always returns a value of opposite parity of the parameter.
\begin{verbatim}
    int hash_it(int x) {
        return x*x+6*x+9;
    }
\end{verbatim}
\end{pro}

\begin{pro}\label{pro:ab1}
Prove or disprove that the following method computes the absolute value of \verb+x+.
For simplicity, assume that all of the calculations are performed with perfect precision.
You may use the fact that $\sqrt{x^2}=x$ when $x\geq 0$ if it will help.
\begin{lstlisting}
    double absoluteValue(double x) {
        double square = x*x;
        double answer = sqrt(square);
        return answer;
    }
\end{lstlisting}
\end{pro}


\begin{pro}\label{pro:ab2}
Prove or disprove that the following method computes the absolute value of \verb+x+.
For simplicity, assume that all of the calculations are performed with perfect precision.
You may use the fact that $(\sqrt{x})^2=x$ when $x\geq 0$ if it will help.
\begin{code}
    double absoluteValue(double x) {
        double root = sqrt(x);
        double answer = root*root;
        return answer;
    }
\end{code}
\end{pro}
\begin{pro}
Problems~\ref{pro:ab1} and~\ref{pro:ab2} both assumed that 
``all of the calculations are performed with perfect precision''.  Is that a realistic assumption?
Give an example of an input for which the each algorithm will work properly.  
Then give an example of an input for which each algorithm will {\em not} work properly.
You can implement and run the algorithms to do some testing if you wish.
\end{pro}

\begin{pro}
The following method is supposed to do some computations on a positive number that result in getting the 
original number back.
Prove or disprove that this method always returns the {\em exact} value that was passed in.
Unlike in the previous problems, here you should assume that although a 
\verb+double+ stores a real number as accurately as possible, 
it uses only a fixed amount of space.  Thus, a \verb+double+ is unable to store the exact value of any
irrational number--it instead stores an approximation.
\begin{code}
    double returnTheParameterUnmodified(double x) {
        double a = sqrt(x);
        double b = a*a;
        return b;
    }
\end{code}
\end{pro}

\begin{pro}\label{pro:twoVarSwap}
Prove or disprove that the algorithm from Example~\ref{exa:swapnothird} actually {\em does} work\footnote{When we say ``works,'' we mean
for {\em all} possible values of $x$ and $y$.}
properly with integer data types stored using 2's 
complement.\footnote{We assume you have previously encountered the
2's complement representation of integers. 
If not, do an Internet search for details.}
You may restrict to 8-bit numbers
if it will help you think about it more clearly--a proof/counterexample for 8-bit number can easily
be modified to work for 32- or 64-bit numbers.
(Hint:  If it doesn't work, what sort of numbers might it fail on?)
\end{pro}

\begin{pro}
Let \verb+A+ and \verb+B+ be \verb+TreeSet+s (See Example~\ref{exa:TreeSet}).
\begin{lenum}
\item  
The method \verb+addAll(TreeSet other)+
{\em adds all of the elements in \verb+other+ to this set if they're not already present.}
What is the result of \verb+A.addAll(B)+ (in terms of \verb+A+ and \verb+B+ and set operators)?
\item 
The method \verb+removeAll(TreeSet other)+
{\em removes from this set all of its elements that are contained in \verb+other+.}
What is the result of \verb+A.removeAll(B)+ (in terms of \verb+A+ and \verb+B+ and set operators)?
\item
Write \verb+A.contains(x)+ using set notation, where $x$ is an element that can be stored in a \verb+TreeSet+.
\end{lenum}
\end{pro}


\begin{pro}
The class \scode{Relation} is a partial implementation of a relation on a set $A$.
It has a list of \scode{Element} objects.  
\begin{itemize}
\item
An \scode{Element} stores an ordered pair from $A$.  
\scode{Element} has methods \scode{getFrom()} and \scode{getTo()} 
(using the language of the directed graph representation).  
So if an \scode{Element} is storing $(a,b)$, \scode{getFrom()} returns $a$ and \scode{getTo()} returns $b$.  
The constructor \scode{Element(Object a, Object b)} creates an element $(a,b)$.
\item The 
\scode{Relation} class has methods like \scode{areRelated(Object a,Object b)}, 
\scode{getElements( )},  and \scode{getUniverse( )}.
\item
Methods in the \scode{Relation} class can use \scode{for(Element e : getElements())} to iterate over
elements of the relation.   
\item
Similarly, the loop \scode{for(Object a : getUniverse())}  iterates over the elements of $A$.
\end{itemize}
Given all of this, implement the following methods in the \scode{Relation} class: 
\begin{lenum}
\item
\scode{isReflexive()}
\item
\scode{isSymmetric()}
\item
\scode{isAntiSymmetric()}
\end{lenum}
\end{pro}
